\chapter{Teoretická část studentské práce}
Tato kapitola shrnuje technologie, na kterých je navržená chytrá domácnost
postavena. Nejprve je popsán standard IEEE~802.15.4, který definuje fyzickou
vrstvu a vrstvu řízení přístupu k médiu pro nízkopříkonové bezdrátové sítě.
Na něm staví síť Thread, které je věnována další sekce. Následuje popis
aplikačního protokolu Matter, a to včetně organizace Connectivity Standards
Alliance a variant přenosu přes Wi-Fi a Thread. Dále je představen nástroj
ESPHome a nejrozšířenější ekosystémy chytré domácnosti, tedy Home Assistant,
Google Home, Apple Home a Samsung SmartThings. Kapitola je zakončena popisem
role OpenThread Border Router, který propojuje síť Thread s ostatními
sítěmi.

\section{IEEE 802.15.4}
Standard IEEE 802.15.4 definuje fyzickou vrstvu (PHY) a vrstvu řízení přístupu k médiu (MAC) pro bezdrátové osobní sítě s nízkou přenosovou rychlostí a velmi nízkou spotřebou~\cite{ieee-802154}. Neurčuje však vyšší vrstvy sítě. Stává se tak základem pro protokoly, které na něm staví, například ZigBee, WirelessHART nebo 6LoWPAN~\cite{mathworks-802154}. Protokol Thread využívá právě fyzickou a MAC vrstvu tohoto standardu~\cite{espressif-openthread}. V pásmu 2,4~GHz, které je používáno ve světě nejčastěji, standard definuje 16~kanálů s odstupem 5~MHz a přenosovou rychlost 250~kb/s~\cite{jain-802154}. Maximální velikost rámce MAC je 127~bajtů~\cite{jain-802154}. Zařízení mohou používat dlouhou 64bitovou adresu nebo krátkou 16bitovou adresu~\cite{nsnam-lrwpan}. Přístup ke sdílenému kanálu řídí mechanismus CSMA/CA, ve kterém zařízení před vysíláním nejprve vyhodnotí, zda je kanál volný, a při obsazeném kanálu vyčká náhodnou dobu~\cite{nxp-an1035}. Standard rozlišuje plnohodnotná zařízení (FFD), která mohou pracovat jako koordinátor, a zařízení s omezenou funkcí (RFD)~\cite{jain-802154}. Toto rozdělení se promítá i do rolí v síti Thread. Čip ESP32-C6 obsahuje vestavěné rádio IEEE~802.15.4, takže může fungovat jako uzel sítě Thread~\cite{espressif-openthread}. 



\section{Thread}



\section{Matter}



\subsection{Connectivity Standards Alliance}
Connectivity Standards Alliance (CSA) je mezinárodní sdružení technologických
společností, které vytváří a udržuje otevřené standardy pro internet věcí.
Do roku 2021 neslo název Zigbee Alliance a jeho vznik je datován do
roku~2002~\cite{csa-matter-arrives}. K přejmenování na CSA došlo
11.~května 2021, kdy bylo zároveň oznámeno nové označení Matter pro
projekt Connected Home over IP (CHIP). Sdružení si ponechalo značku
Zigbee a nadále ji rozvíjí~\cite{csa-rebrand}.

Specifikace Matter 1.0 byla vydána 4.~října 2022 spolu s otevřením
certifikačního programu. V době vydání sdružovala CSA více než
550~společností, z nichž přes 280 se podílelo na vzniku
specifikace~\cite{csa-matter-arrives}. Součástí certifikace jsou testovací
nástroje a akreditované testovací laboratoře, které ověřují nejen samotný
Matter, ale i podkladové sítě Wi-Fi a Thread~\cite{csa-matter-arrives}.

Výrobce, který chce produkt s podporou Matter certifikovat, musí být členem
CSA a získat od ní identifikátor výrobce (Vendor ID), jenž se vyskytuje ve
firmwaru zařízení a v atestačních certifikátech~\cite{matter-handbook}.


\subsection{Matter over Wi-Fi}



\subsection{Matter over Thread}



\section{ESPHome}


\section{Ekosystémy}


\subsection{Home Assistant}


\subsection{Google Home}


\subsection{Apple Home}


\subsection{Samsung SmartThings}


\section{Open Thread Border Router}


